\documentclass[12pt,a4paper]{amsart}

\usepackage{amsmath,amssymb,amsthm}
\usepackage{mathtools}
\usepackage[utf8]{inputenc}
\usepackage[T1]{fontenc}
\usepackage{hyperref}
\usepackage{enumitem,booktabs,array}
\usepackage{geometry}
\geometry{margin=2.5cm}

% -------------------------------------------------------
% Theorem-Umgebungen
% -------------------------------------------------------
\newtheorem{theorem}{Satz}[section]
\newtheorem{lemma}[theorem]{Lemma}
\newtheorem{corollary}[theorem]{Korollar}
\newtheorem{proposition}[theorem]{Proposition}
\theoremstyle{definition}
\newtheorem{definition}[theorem]{Definition}
\newtheorem{example}[theorem]{Beispiel}
\newtheorem{remark}[theorem]{Bemerkung}
\newtheorem{conjecture}[theorem]{Vermutung}

% -------------------------------------------------------
% Eigene Befehle
% -------------------------------------------------------
\newcommand{\N}{\mathbb{N}}
\newcommand{\Z}{\mathbb{Z}}
\newcommand{\Q}{\mathbb{Q}}
\newcommand{\R}{\mathbb{R}}
\newcommand{\C}{\mathbb{C}}
\newcommand{\F}{\mathbb{F}}
\newcommand{\Zp}{\mathbb{Z}_p}
\newcommand{\Qp}{\mathbb{Q}_p}
\newcommand{\OK}{\mathcal{O}_K}
\newcommand{\p}{\mathfrak{p}}
\newcommand{\Lam}{\Lambda}
\DeclareMathOperator{\Gal}{Gal}
\DeclareMathOperator{\Cl}{Cl}
\DeclareMathOperator{\ord}{ord}
\DeclareMathOperator{\rank}{rang}
\DeclareMathOperator{\Ann}{Ann}
\newcommand{\cylp}{\mathbb{Q}(\mu_{p^\infty})}

% -------------------------------------------------------
\begin{document}
% -------------------------------------------------------

\title{Iwasawa-Theorie:\\
       $p$-adische $L$-Funktionen, die Iwasawa-Algebra\\
       und die Hauptvermutung}

\author{Michael Fuhrmann}
\address{Specialist Mathematics Project, Build 119}
\email{specialist-maths@localhost}
\date{\today}

\begin{abstract}
Die Iwasawa-Theorie untersucht das Verhalten arithmetischer Objekte
(Klassengruppen, Selmer-Gruppen, Galois-Kohomologie), wenn man einen
unendlichen Turm von Zahlkörpern aufsteigt.  Die zentralen Objekte sind
die Iwasawa-Algebra $\Lambda = \Zp\llbracket T \rrbracket$ (ein vollständiger
lokaler Ring), $p$-adische $L$-Funktionen (analytisch konstruiert nach
Kubota--Leopoldt) und Iwasawa-Moduln (endlich erzeugte $\Lambda$-Moduln,
charakterisiert durch $\mu$- und $\lambda$-Invarianten).
Wir formulieren die Hauptvermutung (bewiesen von Mazur--Wiles 1984 und
Wiles 1990) und skizzieren die Beweisideen.
Vorausgesetzt werden Kenntnisse der algebraischen Zahlentheorie
(Ganzheitsringe, Klassengruppen, Galois-Theorie).
\end{abstract}

\maketitle
\tableofcontents

% ================================================================
\section{Einleitung und Motivation}
% ================================================================

Sei $p$ eine ungerade Primzahl.  Kummers Kriterium (1850) besagt:
$p$ teilt die Klassenzahl $h(\Q(\zeta_p))$ des $p$-ten Kreisteilungskörpers
genau dann, wenn $p$ \emph{irregulär} ist, d.\,h.\ $p$ teilt den Zähler
einer Bernoulli-Zahl $B_{2k}$ für ein $k = 1, \ldots, (p-3)/2$.

Iwasawas Idee (um 1959) war es, nicht einen einzelnen Kreisteilungskörper,
sondern den gesamten \emph{Kreisteilungsturm} zu betrachten:
\[
  \Q \subset \Q(\zeta_p) \subset \Q(\zeta_{p^2}) \subset \cdots
  \subset \Q(\zeta_{p^\infty}) = \cylp.
\]
Setzt man $K_n = \Q(\zeta_{p^n})$, $K_\infty = \bigcup_n K_n$ und
$\Gamma = \Gal(K_\infty/\Q(\zeta_p)) \cong \Zp$ (eine pro-$p$-Gruppe
isomorph zu den $p$-adischen ganzen Zahlen), so untersucht man den
$p$-Anteil der Klassengruppe $A_n$ von $K_n$ für wachsendes $n$.

\begin{theorem}[Iwasawa-Wachstumsformel]
\label{thm:wachstum}
Es gibt ganze Zahlen $\mu \ge 0$, $\lambda \ge 0$ und $\nu$ (alle abhängig
von $K_\infty/\Q$), sodass für alle hinreichend großen $n$:
\[
  v_p(|A_n|) = \mu p^n + \lambda n + \nu,
\]
wobei $v_p$ die $p$-adische Bewertung ist.
\end{theorem}

Das Wachstum der $p$-adischen Klassenzahlen wird also durch genau zwei
ganze Zahlen $\mu$ und $\lambda$ (die \emph{Iwasawa-Invarianten}) gesteuert.

% ================================================================
\section{Die Iwasawa-Algebra}
% ================================================================

\begin{definition}[Iwasawa-Algebra]
\label{def:iwasawa_algebra}
Sei $\Gamma \cong \Zp$ eine topologische Gruppe (z.\,B.\ die Galois-Gruppe
der zyklotomischen $\Zp$-Erweiterung).  Die \emph{Iwasawa-Algebra} ist der
vervollständigte Gruppenring:
\[
  \Lambda = \Zp\llbracket\Gamma\rrbracket
  = \varprojlim_n \Zp[\Gamma/\Gamma^{p^n}]
  \cong \Zp\llbracket T \rrbracket,
\]
der \emph{formale Potenzreihenring} in einer Variablen über $\Zp$, wobei
$T = \gamma - 1$ für einen topologischen Erzeuger $\gamma$ von $\Gamma$.
\end{definition}

\begin{theorem}[Struktur von $\Lambda$]
\label{thm:lambda_struktur}
Die Iwasawa-Algebra $\Lambda = \Zp\llbracket T \rrbracket$ ist:
\begin{enumerate}[label=(\roman*)]
  \item ein \emph{vollständiger lokaler Ring} mit maximalem Ideal $(p, T)$,
  \item ein \emph{regulärer lokaler Ring} der (Krull-)Dimension $2$,
  \item ein \emph{faktorieller Ring} (UFD),
  \item \emph{noethersch}.
\end{enumerate}
\end{theorem}

\begin{theorem}[Weierstraßscher Vorbereitungssatz für $\Lambda$]
\label{thm:weierstrass}
Jedes von Null verschiedene $f \in \Lambda = \Zp\llbracket T \rrbracket$
lässt sich eindeutig schreiben als
\[
  f(T) = p^\mu \cdot u(T) \cdot P(T),
\]
wobei $\mu \ge 0$, $u(T) \in \Lambda^\times$ eine Einheit und
$P(T) \in \Zp[T]$ ein \emph{ausgezeichnetes Polynom} (normiert vom Grad
$\lambda \ge 0$ mit allen nicht-führenden Koeffizienten durch $p$ teilbar)
ist.  Die ganzen Zahlen $\mu$ und $\lambda$ heißen
\emph{Iwasawa-$\mu$- und $\lambda$-Invariante} von $f$.
\end{theorem}

\begin{remark}
Die Invariante $\mu(f)$ und $\lambda(f) = \deg P$ messen, wie ``tief''
$f$ im maximalen Ideal liegt.  Für das charakteristische Ideal eines
Iwasawa-Moduls steuern sie das Wachstum in Satz~\ref{thm:wachstum}.
\end{remark}

% ================================================================
\section{Endlich erzeugte $\Lambda$-Moduln}
% ================================================================

\begin{theorem}[Struktursatz für $\Lambda$-Moduln]
\label{thm:struktur}
Jeder endlich erzeugte Torsions-$\Lambda$-Modul $M$ ist pseudo-isomorph
zu einer direkten Summe:
\[
  M \;\sim\; \bigoplus_{i=1}^s \Lambda/(p^{\mu_i}) \;\oplus\;
             \bigoplus_{j=1}^t \Lambda/(f_j(T)),
\]
wobei $f_j$ irreduzible ausgezeichnete Polynome in $\Zp[T]$ sind.
\emph{Pseudo-isomorph} bedeutet: Es gibt einen $\Lambda$-Modulhomomorphismus
mit endlichem Kern und Kokern.
\end{theorem}

\begin{definition}[Charakteristisches Ideal und Iwasawa-Invarianten]
\label{def:char_ideal}
Für einen endlich erzeugten Torsions-$\Lambda$-Modul $M$ ist das
\emph{charakteristische Ideal}
\[
  \mathrm{char}(M) = \Bigl(p^{\mu} \cdot \prod_j f_j(T)\Bigr) \subseteq \Lambda,
\]
mit $\mu = \sum_i \mu_i$.
Die \emph{$\mu$-Invariante} $\mu(M) = \sum_i \mu_i$ und die
\emph{$\lambda$-Invariante} $\lambda(M) = \sum_j \deg f_j$ sind
wohldefiniert.
\end{definition}

% ================================================================
\section{Die zyklotomische $\Zp$-Erweiterung}
% ================================================================

\begin{definition}[Zyklotomische $\Zp$-Erweiterung]
\label{def:zyklotom_erw}
Sei $F$ ein Zahlkörper.  Die \emph{zyklotomische $\Zp$-Erweiterung}
von $F$ ist $F_\infty = F(\mu_{p^\infty})^{\mathrm{conn}}$, die eindeutige
$\Zp$-Erweiterung von $F$ in $F(\mu_{p^\infty})$.  Für $F = \Q$:
\[
  \Q_\infty = \bigcup_{n \ge 1} \Q(\zeta_{p^n})^+
  \quad\text{(der reelle Teilkörper von $\Q(\zeta_{p^n})$)}.
\]
Die Galois-Gruppe $\Gamma = \Gal(\Q_\infty/\Q) \cong \Zp$.
\end{definition}

\begin{definition}[Der Iwasawa-Modul $X$]
Sei $M_n$ die maximale abelsche $p$-Erweiterung von $K_n = \Q(\zeta_{p^n})$,
die außerhalb von $p$ unverzweigt ist.  Definiere
\[
  X = \Gal(M_\infty / K_\infty)
  = \varprojlim_n \Gal(M_n/K_n).
\]
Dann ist $X$ ein endlich erzeugter Torsions-$\Lambda$-Modul.
\end{definition}

% ================================================================
\section{Die Kubota--Leopoldt $p$-adische $L$-Funktion}
% ================================================================

\begin{definition}[Spezielle Werte Dirichletscher $L$-Funktionen]
Für einen Dirichlet-Charakter $\chi$ mit Führer $f$ und $\chi(-1) = -1$
(ungerader Charakter) sind die Werte $L(1-n, \chi)$ für $n = 1, 2, 3, \ldots$
algebraische Zahlen:
\[
  L(1-n, \chi) = -\frac{B_{n,\chi}}{n},
  \qquad
  B_{n,\chi} = f^{n-1} \sum_{a=1}^{f} \chi(a) B_n\!\Bigl(\frac{a}{f}\Bigr),
\]
wobei $B_n(x)$ die Bernoulli-Polynome sind.
\end{definition}

\begin{theorem}[Kubota--Leopoldt, 1964]
\label{thm:kubota_leopoldt}
Sei $\chi$ ein primitiver Dirichlet-Charakter mit Führer $f$ und $\chi(-1) = -1$.
Sei $p$ eine ungerade Primzahl mit $p \nmid f$.
Es gibt eine eindeutige stetige Funktion
\[
  L_p(s, \chi) : \Zp \longrightarrow \Qp,
\]
die \emph{Kubota--Leopoldt $p$-adische $L$-Funktion}, welche
\[
  L_p(1-n, \chi) = \bigl(1 - \chi(p)\omega^{-n}(p) p^{n-1}\bigr)
                   L(1-n, \chi\omega^{-n})
\]
für alle positiven ganzen Zahlen $n \geq 1$ erfüllt,
wobei $\omega$ der Teichmüller-Charakter ist.
\end{theorem}

\begin{proof}[Konstruktionsskizze]
\textbf{Schritt 1: Das $p$-adische Maß.}
Man konstruiert ein \emph{$p$-adisches Maß} $\mu_\chi$ auf $\Zp^\times$
(eine beschränkte, finit-additive $\Qp$-wertige Funktion auf offenen
kompakten Teilmengen) mit
\[
  \int_{\Zp^\times} x^{n-1} \, d\mu_\chi(x)
  = (1 - \chi(p)\omega^{-n}(p)p^{n-1}) L(1-n, \chi\omega^{-n}).
\]

\textbf{Schritt 2: $p$-adische Interpolation.}
Die Werte $(1 - \chi(p)p^{n-1}) L(1-n, \chi)$ für $n \equiv 0 \pmod{p-1}$
sind $p$-adisch nahe für $n$ in arithmetischen Progressionen.
Es gibt eine Potenzreihe $G_\chi(T) \in \Zp\llbracket T \rrbracket$
mit $G_\chi((1+p)^{1-n} - 1) = $ den interpolierten Werten.

\textbf{Schritt 3: Iwasawa-Potenzreihe.}
Über die Identifikation $\Gamma \cong \Zp$ (wobei ein topologischer
Erzeuger $\gamma$ auf $1+p$ abgebildet wird) erhält man
\[
  L_p(s, \chi) = G_\chi((1+p)^s - 1) \in \Lambda \otimes_{\Zp} \Qp.
\]
\end{proof}

\begin{remark}[Triviale Nullstelle]
Für $\chi = \omega^k$ kann $L_p(s, \omega^k)$ bei $s = 1-k$ verschwinden,
auch wenn $L(1-k, \omega^k) \ne 0$.  Dieses Phänomen der
\emph{trivialen Nullstelle} wurde von Ferrero--Greenberg und
später von Greenberg--Stevens untersucht.
\end{remark}

% ================================================================
\section{Die Hauptvermutung (Satz von Mazur--Wiles)}
% ================================================================

\begin{conjecture}[Iwasawas Hauptvermutung, ursprüngliche Form]
\label{conj:haupt}
Sei $p$ eine ungerade Primzahl.  Für jede gerade ganze Zahl $i$
mit $0 \le i \le p-3$ sei $X^{(i)}$ der $\omega^i$-Eigenraum von $X$
und $f_i(T) \in \Lambda$ ein Erzeuger des charakteristischen Ideals
$\mathrm{char}(X^{(i)})$.  Dann gilt:
\[
  f_i(T) \;\sim\; G_{\omega^{1-i}}(T) \pmod{\text{Einheiten}},
\]
wobei $G_{\omega^{1-i}}(T)$ die Iwasawa-Potenzreihe der $p$-adischen
$L$-Funktion zum Charakter $\omega^{1-i}$ ist.
\end{conjecture}

\begin{theorem}[Mazur--Wiles 1984; Wiles 1990]
\label{thm:mazur_wiles}
Die Iwasawa-Hauptvermutung gilt für alle ungeraden Primzahlen $p$ und
alle geraden Charaktere $\omega^i$ wie oben.
\end{theorem}

\begin{proof}[Beweissskizze (Wiles' Ansatz)]
Wiles' Beweis von 1990 verläuft in drei Hauptschritten:

\textbf{Schritt 1: Konstruktion eines Euler-Systems.}
Man konstruiert \emph{Euler-System-Elemente} $c_n \in K_n^\times$
(bezogen auf Kreiseinheiten und Gauß-Summen), deren Normen kompatibel
durch den Turm sind.  Die Hauptquelle sind die \emph{Kreiseinheiten}:
$\eta_n = $ das Bild von $(1 - \zeta_{p^n})$ in $K_n^\times$.

\textbf{Schritt 2: Indexformel für Kreiseinheiten.}
Sei $C_n$ die Gruppe der Kreiseinheiten in $K_n$.  Es gilt
\[
  [E_n : C_n] = h_n^+ \quad\text{(die reelle Klassenzahl)},
\]
bewiesen mit der analytischen Klassenzahlformel.

\textbf{Schritt 3: Hecke-Algebren, Modulkurven und Vergleich der Fitting-Ideale.}
Wiles (1990) etablierte die zentrale Teilbarkeit
\[
  \mathrm{char}(X^{(i)}) \mid (G_{\omega^{1-i}}(T)) \quad \text{in } \Lambda
\]
mithilfe von \emph{Hecke-Algebren und Modulkurven}, in Erweiterung der Strategie
von Mazur--Wiles (1984): Man verbindet den Iwasawa-Modul $X^{(i)}$ mit der
$p$-adischen Hecke-Algebra, die auf der Jacobischen der Modulkurve $X_1(p^n)$
wirkt, im Anschluss an Ribets Beweis des Umkehrssatzes von Herbrand.
Die Gleichheit der charakteristischen Ideale folgt durch Vergleich der
$\mu$-Invarianten: Ferrero--Washington (1979) bewiesen $\mu(X^{(i)}) = 0$
für $\Q$-Erweiterungen, und die $\mu$-Invariante von $G_{\omega^{1-i}}$
ist ebenfalls $0$ (wegen $L(1, \omega^{1-i}) \ne 0$).  Gleiche $\mu$-Invarianten
erzwingen durch ein Rangargument die Gleichheit der charakteristischen Ideale.

\textit{Hinweis:} Der Euler-System-Ansatz (direkte Verwendung normkompatibler
Kreiseinheiten) wurde von Rubin (1991) für imaginär-quadratische Körper entwickelt
und liefert einen eleganteren Beweis, ist aber nicht die von Wiles 1990 verwendete Methode.
\end{proof}

\begin{remark}[Spätere Verallgemeinerungen]
\begin{itemize}
  \item Rubin (1991) gab einen kürzeren Beweis mit Euler-Systemen für
    imaginär-quadratische Körper.
  \item Kato (2004) bewies eine Teilbarkeit der Hauptvermutung für
    elliptische Kurven über $\Q$ mittels Beilinson--Kato-Elementen.
  \item Die \emph{nicht-kommutative} Hauptvermutung für $p$-adische
    Darstellungen wurde von Coates--Fukaya--Kato--Sujatha--Venjakob (2005)
    formuliert.
\end{itemize}
\end{remark}

% ================================================================
\section{Die $\mu$- und $\lambda$-Invarianten}
% ================================================================

\begin{definition}[$\mu$- und $\lambda$-Invarianten eines Zahlkörpers]
\label{def:mu_lambda}
Für die zyklotomische $\Zp$-Erweiterung $K_\infty/K$ sind die
$\mu$- und $\lambda$-Invarianten diejenigen des Iwasawa-Moduls $X$:
\[
  \mu(K/p) = \mu(X), \qquad \lambda(K/p) = \lambda(X).
\]
Nach Satz~\ref{thm:wachstum} gilt für große $n$:
$v_p(|A_n|) = \mu p^n + \lambda n + \nu$.
\end{definition}

\begin{conjecture}[Iwasawas $\mu = 0$-Vermutung]
Für jeden Zahlkörper $K$ und jede Primzahl $p$ gilt
$\mu(K_\infty/K) = 0$.
\end{conjecture}

\begin{theorem}[Ferrero--Washington, 1979]
Für $K = \Q$ und jede ungerade Primzahl $p$ gilt $\mu(\Q_\infty/\Q) = 0$.
Allgemeiner: für jede abelsche Erweiterung $K/\Q$ gilt $\mu(K_\infty/K) = 0$.
\end{theorem}

\begin{proof}[Beweisidee]
Ferrero und Washington zeigten, dass die $p$-adische $L$-Funktion
$L_p(s, \chi)$ nicht identisch $\equiv 0 \pmod{p}$ ist, was $\mu = 0$
erzwingt.  Ihr Argument nutzt die Theorie $p$-adischer Maße und die
explizite Form der Bernoulli-Zahlen.
\end{proof}

% ================================================================
\section{BSD und Iwasawa-Theorie für elliptische Kurven}
% ================================================================

\begin{definition}[Selmer-Gruppe und Iwasawa-Modul für $E/\Q$]
\label{def:selmer_iwasawa}
Sei $E/\Q$ eine elliptische Kurve und $p$ eine ungerade Primzahl mit
gewöhnlicher guter Reduktion ($a_p(E) \not\equiv 0 \pmod{p}$).
Die \emph{$p$-adische Selmer-Gruppe} über $\Q_\infty$ ist
\[
  \mathrm{Sel}_{p^\infty}(E/\Q_\infty)
  = \ker\Bigl(H^1(\Q_\infty, E[p^\infty]) \to
    \prod_v H^1(\Q_{\infty,v}, E)[p^\infty]/E(\Q_{\infty,v})\Bigr).
\]
Der Pontryagin-Dual $X_E = \mathrm{Sel}_{p^\infty}(E/\Q_\infty)^\vee$
ist ein endlich erzeugter $\Lambda$-Modul.
\end{definition}

\begin{conjecture}[Iwasawa-Hauptvermutung für elliptische Kurven]
\label{conj:ec_haupt}
Für $E/\Q$ mit gewöhnlicher guter Reduktion bei $p$ gibt es eine
$p$-adische $L$-Funktion $L_p(E, s) \in \Lambda$
(die Mazur--Swinnerton-Dyer $p$-adische $L$-Funktion) mit
\[
  \mathrm{char}(X_E) = (L_p(E, s)) \quad \text{in } \Lambda.
\]
\end{conjecture}

\begin{theorem}[Kato 2004; Skinner--Urban 2014]
\label{thm:kato_su}
Eine Teilbarkeit $\mathrm{char}(X_E) \mid (L_p(E, s))$ gilt unter milden
Voraussetzungen (Kato 2004).  Unter weiteren Bedingungen gilt die
Hauptvermutung (Skinner--Urban 2014) und insbesondere:
\[
  L_p(E, 1) = 0 \implies \mathrm{rang}(E(\Q)) \ge 1.
\]
\end{theorem}

\begin{remark}[Verbindung zu BSD]
Die Iwasawa-Hauptvermutung für $E$ impliziert den $p$-Anteil der
Birch--Swinnerton-Dyer-Vermutung.  Die vollständige BSD-Vermutung
ist weiterhin offen.
\end{remark}

% ================================================================
\section{Explizite Beispiele}
% ================================================================

\begin{example}[$p=3$, $\Q(\zeta_3)$]
$\zeta_3 = e^{2\pi i/3}$.
$h(\Q(\zeta_3)) = h(\Q(\zeta_9)) = h(\Q(\zeta_{27})) = 1$.
Durch die Wachstumsformel folgt $\mu = \lambda = 0$.
(Die kleinste irreguläre Primzahl, die $3$ teilt, ist nicht in diesem Turm.)
\end{example}

\begin{example}[Irreguläre Primzahl $p = 37$]
$p = 37$ ist die kleinste irreguläre Primzahl: $37 \mid B_{32}$.
Die Klassenzahl $h(\Q(\zeta_{37})) = 37$.
Die Iwasawa-$\lambda$-Invariante erfüllt $\lambda \ge 1$.
Berechnungen (Buhler--Crandall 1992) ergeben $\mu = 0$, $\lambda = 2$.
\end{example}

% ================================================================
\section{Zusammenfassung und Ausblick}
% ================================================================

\begin{remark}[Was die Iwasawa-Theorie leistet]
Die Hauptvermutung (jetzt Satz) etabliert ein grundlegendes Wörterbuch:
\begin{center}
\begin{tabular}{l|l}
\textbf{Analytische Seite} & \textbf{Algebraische Seite} \\
\hline
$p$-adische $L$-Funktion $L_p(s, \chi)$ & Char.\ Ideal von $X^{(i)}$ \\
$L(1-n, \chi) \ne 0$ & $X^{(i)}$ ist endlich \\
$\mu(G_\chi) = 0$ (Ferrero--Washington) & $\mu(X^{(i)}) = 0$ \\
$\lambda(G_\chi)$ & Rang von $X^{(i)}$ über $\Lambda$ \\
\end{tabular}
\end{center}
\end{remark}

\begin{remark}[Offene Probleme]
\begin{itemize}
  \item \textbf{$\mu = 0$-Vermutung} für nicht-abelsche Erweiterungen.
  \item \textbf{Nicht-kommutative Hauptvermutung} für GL$_n$-Motive.
  \item \textbf{$p$-adische BSD}: Vollständige BSD aus der Iwasawa-Theorie
    erfordert die Kontrolle der Tate--Shafarevich-Gruppe.
  \item \textbf{$p$-adische Langlands-Korrespondenz}: Das Breuil--Colmez-Programm
    zielt darauf ab, Iwasawa-Moduln via $(\varphi,\Gamma)$-Moduln zu kodieren.
\end{itemize}
\end{remark}

\section*{Danksagung}
Der Autor dankt dem Specialist Mathematics Project für die Bereitstellung
der Forschungsinfrastruktur.

\begin{thebibliography}{99}

\bibitem{Iwasawa1959}
K.~Iwasawa,
\textit{On $\Gamma$-extensions of algebraic number fields},
Bull.\ Amer.\ Math.\ Soc.\ \textbf{65} (1959), 183--226.

\bibitem{KubotaLeopoldt1964}
T.~Kubota und H.-W.~Leopoldt,
\textit{Eine $p$-adische Theorie der Zetawerte, I: Einführung der $p$-adischen Dirichletschen $L$-Funktionen},
J.\ Reine Angew.\ Math.\ \textbf{214/215} (1964), 328--339.

\bibitem{MazurWiles1984}
B.~Mazur und A.~Wiles,
\textit{Class fields of abelian extensions of $\Q$},
Invent.\ Math.\ \textbf{76} (1984), 179--330.

\bibitem{Wiles1990}
A.~Wiles,
\textit{The Iwasawa conjecture for totally real fields},
Ann.\ of Math.\ (2) \textbf{131} (1990), 493--540.

\bibitem{Rubin1991}
K.~Rubin,
\textit{The ``main conjecture'' of Iwasawa theory for imaginary quadratic fields},
Invent.\ Math.\ \textbf{103} (1991), 25--68.

\bibitem{FerWash1979}
B.~Ferrero und L.~C.~Washington,
\textit{The Iwasawa invariant $\mu_p$ vanishes for abelian number fields},
Ann.\ of Math.\ (2) \textbf{109} (1979), 377--395.

\bibitem{Kato2004}
K.~Kato,
\textit{$p$-adic Hodge theory and values of zeta functions of modular forms},
Ast\'erisque \textbf{295} (2004), 117--290.

\bibitem{SkinnerUrban2014}
C.~Skinner und E.~Urban,
\textit{The Iwasawa main conjecture for $\mathrm{GL}_2$},
Invent.\ Math.\ \textbf{195} (2014), 1--277.

\bibitem{Neukirch1999}
J.~Neukirch,
\textit{Algebraische Zahlentheorie},
Springer, Berlin, 1992.

\bibitem{Washington1997}
L.~C.~Washington,
\textit{Introduction to Cyclotomic Fields},
2.\ Aufl., Springer, New York, 1997.

\bibitem{Coates2006}
J.~Coates und R.~Sujatha,
\textit{Cyclotomic Fields and Zeta Values},
Springer, Berlin, 2006.

\end{thebibliography}

\end{document}
